\PilotPage{6}
Define the natural frequency in radians per second as
\begin{equation}
    \omega_n=\sqrt{\frac{g}{L}}.
    \tag{10}
\end{equation}
The solutions are
\begin{equation}
\begin{aligned}
    s_1 &= i\omega_n,\\
    s_2 &= -i\omega_n.
\end{aligned}
    \tag{11}
\end{equation}
Thus,
\begin{equation}
    \theta(t)=C_1^*e^{i\omega_nt}+C_2^*e^{-i\omega_nt},
    \tag{12}
\end{equation}
or, equivalently,
\begin{equation}
    \theta(t)=A^*\cos\omega_nt+B^*\sin\omega_nt.
    \tag{13}
\end{equation}

The frequency in hertz is
\begin{equation}
    f_n=\frac{\omega_n}{2\pi}=\frac{1}{2\pi}\sqrt{\frac{g}{L}}.
    \tag{14}
\end{equation}
The period of oscillation is
\begin{equation}
    \tau=\frac{1}{f_n}=2\pi\sqrt{\frac{L}{g}}.
    \tag{15}
\end{equation}

\subsection*{Particle motion}

The particle displacement is
\begin{equation*}
    u=L\theta.
\end{equation*}
Consequently,
\begin{align}
    u(t)&=C_1e^{i\omega_nt}+C_2e^{-i\omega_nt},
    \tag{16a}\\
    u(t)&=A\cos\omega_nt+B\sin\omega_nt,
    \tag{16b}\\
    u(t)&=C\sin(\omega_nt+\varphi).
    \tag{16c}
\end{align}
